%! Bias in Samples and Surveys — Unit 1, Lesson 1.4 — Warm-Up (Answer Key)
% PILOT SHAPE (2026-09-06): modeled on AP Statistics lesson 1.4 minus its AP Practice page.
% ONE PAGE, blank and key. Three items at 12pt, \workrowsep 50pt, item spacing 22pt, so each
% item gets real handwriting room and still fits. Do not add a fourth item — the page is full.
% GRADUAL RELEASE: the warm-up rehearses, WITHOUT naming bias, the three moves the notes run
% on. Item 1 is the census truth as a percent (300/1500 = 20%). Item 2 is the Tuesday-evening
% sample as a percent (27/75 = 36%) — the study 1.3 threw out. Item 3 scales 36% to 540 and
% puts it beside the truth: off by 240, too high — the first time in Unit 1 a statistic sits
% beside a KNOWN parameter. Section 1 of the notes opens on these numbers; the hook asks whether
% ten times the sample would fix it.
\documentclass[12pt]{article}
\usepackage{saar-article}
\usepackage{saar-key}   % pulls in -boxes; do NOT also load -boxes
\newcommand{\TallMath}[1]{$\displaystyle #1\rule[-1.4em]{0pt}{3.2em}$}
% Word bank strip for every fill-in cluster (user request 2026-09-06). Defined per document;
% the key inherits it and prints the same strip, so heights match.
\newcommand{\wordbank}[1]{\par\vspace{2pt}\noindent\fcolorbox{goldacc}{goldbg}{\parbox{\dimexpr\linewidth-2\fboxsep-2\fboxrule\relax}{\small\textbf{Word bank:} #1}}\par\vspace{4pt}}
% Handwriting room inside every \begin{work} block. Moves the blank and the
% key together, so the two can never drift apart.
\setlength{\workrowsep}{50pt}

\begin{document}

\pageheader{Unit 1, Lesson 1.4}{Warm-Up --- Answer Key}
% NAMESTRIP: no \namedateperiod here — mirrors the blank. NO teachernote here —
% teacher prose lives in the lesson plan.

\vspace{0.15in}

\begin{notesbox}{Spiral Review --- Harbor Point Community Pool}
\normalsize
Harbor Point Community Pool has $1{,}500$ members. The office keeps a record of what
every member says they mainly come to do.

\begin{enumerate}[leftmargin=1.8em, itemsep=22pt, topsep=10pt]

  \item The office's records show that $300$ of the $1{,}500$ members mainly swim laps.
        What \textbf{percent} of the membership is that?

        \begin{work}
          \dfrac{300}{1500} &= 0.20 \\
                            &= 20\%
        \end{work}

  \item Last class the board threw out this study: a staff member asked the $75$ members
        who arrived one Tuesday evening, and $27$ of them said they mainly swim laps. What
        \textbf{percent} of those $75$ is that?

        \begin{work}
          \dfrac{27}{75} &= 0.36 = 36\%
        \end{work}

  \item The staff member uses her $36\%$ to describe the whole membership.

        \textbf{(a)} How many of the $1{,}500$ members does $36\%$ predict?

        \begin{work}
          0.36 \times 1500 &= 540 \text{ members}
        \end{work}

        \textbf{(b)} The records in item 1 say the real number is $300$. By how many
        members is the prediction off, and is it too high or too low?
        \wordbank{16 \;\textbullet\; 240 \;\textbullet\; 300 \;\textbullet\; too high \;\textbullet\; too low}

        Off by \ans{240} members \qquad The prediction is \ans{too high}

\end{enumerate}
\end{notesbox}

\end{document}
